\section*{Problem 761 (Round 2)}

\subsection*{1) ROUND-2 OBJECTIVE}
I pursue \textbf{(C) obstruction/correction}. A full proof or counterexample is still out of reach, but Round 1 left two decisive gaps that can now be closed rigorously:
\begin{enumerate}
\item the relation between the two questions (large chromatic number versus large cochromatic number), and
\item the lack of any unconditional lower bound on $\delta(G)$ in terms of $\chi(G)$ beyond the trivial $\delta(G)\le \chi(G)$.
\end{enumerate}
The most promising Round-2 step is therefore to isolate the exact structural obstruction by proving the right monotonicity statements, showing that the two formulations are equivalent up to a parameter transfer, and importing the best currently verified quantitative lower bound.

\subsection*{2) Round-1 FOUNDATION USED}
I use the following Round-1 ingredients.
\begin{itemize}
\item For every graph $G$, one has $1\le \delta(G)\le \chi(G)$.
\item The Erd\H{o}s--Neumann-Lara conjecture (Question Q1) was identified correctly as open.
\item The Alon--Krivelevich--Sudakov theorem was already cited in Round 1: large chromatic number forces a subgraph of cochromatic number $\Omega(\chi/\log\chi)$.
\item The fractional result of Mohar--Wu was already cited in Round 1.
\end{itemize}

\subsection*{3) NEW INSIGHT / TOOL (ROUND-2)}
Round 2 adds four concrete advances.
\begin{enumerate}
\item[(i)] \textbf{Subgraph monotonicity of $\delta$.} If $H\subseteq G$, then $\delta(H)\le \delta(G)$.
\item[(ii)] \textbf{Equivalence of the two Erd\H{o}s questions up to parameter transfer.} Since $\zeta(G)\le \chi(G)$, Question~Q1 immediately implies Question~Q2 by taking $H=G$. Conversely, the Alon--Krivelevich--Sudakov theorem converts Question~Q2 back into Question~Q1.
\item[(iii)] \textbf{A genuine quantitative lower bound.} Aboulker, Havet, Pirot and Schabanel proved that every $n$-vertex graph satisfies
\[
\chi(G)\le 2\delta(G)\bigl(1+\lfloor\log_2 n\rfloor\bigr).
\]
Equivalently,
\[
\delta(G)\ge \frac{\chi(G)}{2(1+\lfloor\log_2 |V(G)|\rfloor)}.
\]
This does not settle the conjecture, but it rules out any counterexample family whose order is too small relative to its chromatic number.
\item[(iv)] \textbf{List-analogue evidence.} In 2026, Harutyunyan, Picasarri-Arrieta and Puig i Surroca proved the list analogue
\[
\vec\chi_\ell(G)\ge \Bigl(\frac13-o(1)\Bigr)\log_2 \chi_\ell(G),
\]
which is substantial adjacent evidence even though it does not settle the original parameter $\delta(G)$.
\end{enumerate}

\subsection*{4) ATTACK PLAN (ROUND-2)}
The precise Round-1 gaps were:
\begin{itemize}
\item Round 1 only explained one implication between Q1 and Q2; the reverse implication and the monotonicity needed for it were not made explicit.
\item Round 1 recorded fractional evidence but no unconditional inequality forcing $\delta(G)$ to grow with $\chi(G)$ under an additional hypothesis.
\end{itemize}
I therefore prove the missing monotonicity lemma, deduce the equivalence of Q1 and Q2 (up to a logarithmic loss through the AKS theorem), and then apply the 2025 theorem
\[
\chi(G)\le 2\delta(G)(1+\lfloor\log_2 |V(G)|\rfloor)
\]
to obtain a sharp obstruction to bounded-dichromatic counterexamples.

\subsection*{5) WORK (ROUND-2)}
Recall that $\delta(G)$ denotes the maximum dichromatic number over all orientations of $G$.

\paragraph{Lemma 5.1 (subgraph monotonicity).}
If $H\subseteq G$, then
\[
\delta(H)\le \delta(G).
\]

\emph{Proof.}
Choose an orientation $D_H$ of $H$ with $\chi_d(D_H)=\delta(H)$. Extend $D_H$ arbitrarily to an orientation $D_G$ of $G$. Any dicolouring of $D_G$ restricts to a dicolouring of $D_H$, so
\[
\chi_d(D_H)\le \chi_d(D_G)\le \delta(G).
\]
Therefore $\delta(H)=\chi_d(D_H)\le \delta(G)$. \qed

\paragraph{Proposition 5.2.}
A positive answer to Question~Q1 implies a positive answer to Question~Q2.

\emph{Proof.}
Suppose Q1 holds: there exists a function $f$ such that $\chi(G)\ge f(k)$ implies $\delta(G)\ge k$. Since every graph satisfies $\zeta(G)\le \chi(G)$, the implication
\[
\zeta(G)\ge f(k)\implies \chi(G)\ge f(k)\implies \delta(G)\ge k
\]
shows that Q2 holds, taking $H=G$. \qed

\paragraph{Proposition 5.3.}
A positive answer to Question~Q2 implies a positive answer to Question~Q1.

\emph{Proof.}
Assume Q2 holds, so there exists a function $g$ such that
\[
\zeta(J)\ge g(k)\implies \exists H\subseteq J \text{ with } \delta(H)\ge k.
\]
Let $G$ be a graph with chromatic number $m$. By the Alon--Krivelevich--Sudakov theorem, $G$ contains a subgraph $J\subseteq G$ with
\[
\zeta(J)\ge c\,\frac{m}{\log m}
\]
for some absolute constant $c>0$ and all large enough $m$. Choose $m$ so large that
\[
c\,\frac{m}{\log m}\ge g(k).
\]
Then Q2 gives a subgraph $H\subseteq J\subseteq G$ with $\delta(H)\ge k$. By Lemma 5.1,
\[
\delta(G)\ge \delta(H)\ge k.
\]
Hence large chromatic number forces large dichromatic number, i.e. Q1 holds. \qed

\paragraph{Conclusion from Propositions 5.2 and 5.3.}
The two questions in Problem 761 are qualitatively equivalent: Q1 trivially implies Q2, while Q2 implies Q1 via the AKS subgraph theorem.

\paragraph{External theorem used.}
(Aboulker--Havet--Pirot--Schabanel.) Every graph $G$ of order $n$ satisfies
\[
\chi(G)\le 2\delta(G)\bigl(1+\lfloor\log_2 n\rfloor\bigr).
\]

\paragraph{Corollary 5.4.}
Every graph $G$ satisfies
\[
\delta(G)\ge \frac{\chi(G)}{2(1+\lfloor\log_2 |V(G)|\rfloor)}.
\]
In particular, if
\[
\chi(G)>2(k-1)(1+\lfloor\log_2 |V(G)|\rfloor),
\]
then $\delta(G)\ge k$.

\emph{Proof.}
This is just a rearrangement of the theorem. \qed

\paragraph{Corollary 5.5 (exponential-size obstruction).}
Fix $D\in\mathbb N$. If $\delta(G)\le D$, then
\[
|V(G)|\ge 2^{\chi(G)/(2D)-1}.
\]

\emph{Proof.}
From Corollary 5.4,
\[
\chi(G)\le 2D(1+\lfloor\log_2 |V(G)|\rfloor)\le 2D(1+\log_2 |V(G)|).
\]
Hence
\[
\log_2 |V(G)|\ge \frac{\chi(G)}{2D}-1,
\]
which is equivalent to the displayed bound. \qed

\paragraph{What this means for the original conjecture.}
Any family $\{G_t\}$ with bounded dichromatic number and unbounded chromatic number must necessarily have order at least exponential in $\chi(G_t)$. Thus one entire class of possible counterexamples is ruled out: no counterexample family can have $|V(G_t)|=\exp(o(\chi(G_t)))$.

\subsection*{6) ADVERSARIAL VERIFICATION}
\begin{itemize}
\item \textbf{Monotonicity check.} In Lemma 5.1, extending an orientation from $H$ to $G$ is harmless because every dicolouring of the larger oriented graph restricts to a dicolouring of the smaller one.
\item \textbf{Q1 $\Rightarrow$ Q2 check.} This uses only the elementary inequality $\zeta(G)\le \chi(G)$ and is therefore immediate.
\item \textbf{Q2 $\Rightarrow$ Q1 check.} The only nontrivial input is the AKS theorem. After obtaining $H\subseteq G$ with large $\delta(H)$, Lemma 5.1 is exactly what transfers the conclusion back to $G$.
\item \textbf{Quantifier check.} The equivalence of Q1 and Q2 is only at the level of existence of some forcing function; the actual functions differ because the AKS step loses a logarithm.
\item \textbf{Why the list result does not solve the problem.} Large list dichromatic number does not force equally large ordinary dichromatic number for the same orientation, so the 2026 list theorem is genuine evidence, not a solution to Q1.
\end{itemize}

\subsection*{7) FINAL}
\textbf{UNRESOLVED (BUT STRICTLY ADVANCED).}
Round 2 proves a genuine structural reduction not present in Round 1:
\begin{enumerate}
\item the two questions in Problem 761 are equivalent up to a parameter transfer, and
\item every graph satisfies the unconditional lower bound
\[
\delta(G)\ge \frac{\chi(G)}{2(1+\lfloor\log_2 |V(G)|\rfloor)}.
\]
Thus any bounded-$\delta$ counterexample family must have exponentially large order in its chromatic number. The list analogue is now settled, which further narrows the gap.
\end{enumerate}

\subsection*{8) COMPLETION ESTIMATE (MANDATORY)}
\textbf{COMPLETION: 45\%}

\subsection*{9) REFERENCES}
\begin{itemize}
\item B. Mohar and H. Wu, \emph{Dichromatic number and fractional chromatic number}, Forum Math. Sigma 4 (2016), e32.
\item N. Alon, M. Krivelevich, and B. Sudakov, \emph{Subgraphs with a large cochromatic number}, J. Graph Theory 26 (1997), 295--297.
\item P. Aboulker, F. Havet, F. Pirot, and J. Schabanel, \emph{Minimum acyclic number and maximum dichromatic number of oriented triangle-free graphs of a given order}, Electron. J. Combin. 32(4) (2025), Paper P4.27.
\item A. Harutyunyan, L. Picasarri-Arrieta, and G. Puig i Surroca, \emph{On the list version of a conjecture of Erd\H{o}s and Neumann-Lara}, arXiv:2603.01020, 2026.
\end{itemize}

\subsection*{10) OUTPUT FORMAT}
Here the Erd\H{o}s problem number is $x=761$.

\bigskip
